\begin{recipe}
[% 
    preparationtime = {\unit[2]{h} + \unit[12--24]{h} Ziehzeit},
    bakingtime={\unit[75]{min} + \unit[45]{min} Sauce + \unit[15]{min} Ofen},
    bakingtemperature={\protect\bakingtemperature{
        topbottomheat=\unit[200]{\textcelcius}}},        
    portion = {\portion{2}},
    source = {\protect\recipesource{https://www.instagram.com/p/DUNhVVpiLi7/}{Instagram: vegan.auf.bayrisch}},
    calory = {1050kcal (1755kcal) | 72g (81g) Protein | 78g (211g) KH | 21g (24g) Fett},
    rating = {5},
    price = {4,60€ (6,84€)}
]
{Veganer Sauerbraten (Seitan)}

\graph{% pictures
    small=pic/hauptgericht/03_veganer-sauerbraten-small,
    big=pic/hauptgericht/03_veganer-sauerbraten-big
}

\introduction{%
    Ein veganer Sauerbraten, der perfekt für die Feiertage gemacht ist.
    Kräftiger Sud, lange Ziehzeit und eine intensive, leicht süß-saure Sauce. 
    Nichts für mal eben schnell, aber lohnt sich mega.
}

\ingredients{%
    \textbf{Seitan} & \\
    1000 g & Mehl Typ 550\\
    600 ml & Wasser\\
    1 Prise & Salz\\
    3 TL & Senf\\
    2 TL & Misopaste\\
    2 TL & Sojasauce\\
    1 TL & Paprikapulver\\
    %& Schwarzer Pfeffer\\
    2 TL & Öl\\
    \multicolumn{2}{l}{\textit{\small Fortsetzung auf nächster Seite}}
}
\ingredientsB{%
    \textbf{Sud} & \\
    750 ml & Trockener Rotwein\\
    450 ml & Gemüsebrühe\\
    300 ml & Wasser\\
    5--6 EL & Apfelessig\\
    2 & Große Zwiebeln\\
    2 & Karotten\\
    %80--100 g & Sellerie oder Petersilienwurzel\\
    3 & Lorbeerblätter\\
    6 & Wacholderbeeren\\
    %3 & Nelken\\
    10--12 & Pfefferkörner\\
    2--3 & Getrocknete Shiitakepilze\\
    2 EL & Sojasauce\\
    \textbf{Sauce} & \\
    2 TL & Zucker\\
    2 EL & Tomatenmark\\
    2 TL & Misopaste\\
    2 TL & Preiselbeeren oder Johannisbeergelee\\
    Prise & Salz, Pfeffer\\
    Optional & etwas Essig
}


\preparation{%
    \step%
    \textbf{Teig herstellen:} Mehl, Salz und Wasser zu einem festen Teig kneten. 
    Zu einer Kugel formen und 1 Stunde in kaltem Wasser ruhen lassen.
    
    \step%
    \textbf{Seitan auswaschen:} Teig im kalten Wasser kneten, Wasser mehrfach wechseln, bis es nur noch leicht trüb ist (ca. 10--15 Minuten). 
    Übrig bleibt ca 300gr elastischer Gluten, welchen man 5 Minuten abtropfen lässt.                  
    
    \step%
    \textbf{Würzen \& formen:} Seitan flach drücken, mit Würzpaste bestreichen, in drei Stränge teilen, flechten und zu einer kompakten Bratenform verknoten. 
    Das Ende überstülpen, sodass sich eine Art Haut bildet.
    
    \step%
    \textbf{Sud vorbereiten:} Alle Sud-Zutaten aufkochen und 10--15 Minuten leicht köcheln lassen.
    
    \step%
    \textbf{Garen \& beizen:} Seitan in den heißen (nicht kochenden) Sud legen und bei 85--95 °C 60--75 Minuten ziehen lassen. 
    Im Sud abkühlen lassen und anschließend 12--24 Stunden im Kühlschrank durchziehen lassen.
    
    \step%
    \textbf{Anbraten:} Seitan trocken tupfen und in Öl rundum kräftig anbraten (8--12 Minuten).
    
    \step%
    \textbf{Sauce kochen:} Sud abseihen. Zucker karamellisieren, Tomatenmark rösten, mit Sud ablöschen und auf etwa ein Drittel einkochen (30--45 Minuten). 
    Mit Miso, Preiselbeeren, Salz, Pfeffer und ggf. Essig abschmecken.
    
    \step%
    \textbf{Glasieren:} Seitan mit etwas Sauce in eine Form legen und bei 200 °C etwa 15 Minuten im Ofen glasieren.
}

\suggestion[Beilagen]{%
    Als Beilage empfehle ich Kartoffelklöße und Rotkohl. Fertigen Kloßteig bekommt man auch einfach im Kaufland. Ich habe für die zwei Portionen 750gr Kloßteig und 350gr Rothkohl (auch fertig aus dem Glas) genommen.
}

\suggestion[Verbesserungen]{%
    Die Sauce ist mega lecker und bedarf keiner Verbesserungen, der Seitan an sich war aber eher geschmacksneural, die Haut aber richtig gut. Somit könnte man einerseits beim auslegen und bestreichen das deutlich dünner auslegen, vielleicht sogar die Marinade richtig in den Teig hereinkneten.\\
    Die Struktur vom Seitan nach dem Waschen war sehr gummiartig und nach dem kochen hat es an Hallumi errinert, vielleicht nicht bis zum Ende auswaschen, sondern bis er noch leicht Teigstruktur hat, so könnte auch die Sauce beim einlegen mehr einziehen.
}

\hint{
    Es gitb auch direkt Seitanfix zu kaufen, dabei erspart man sich das nervige auswaschen.
}

\end{recipe}